\chapter{CONCLUSION}
Cette définition de projet de recherche propose une étude qui vise à déterminer la possibilité de représenter fidèlement et exhaustivement une vidéo par une représentation exclusivement textuelle, pour permettre à un grand modèle de langage (LLM) d'identifier automatiquement les moments les plus pertinents de la vidéo. Le projet vise à utiliser les avancées des grands modèles de langages pour analyser une vidéo en utilisant moins de ressources qu'avec une approche multimodale plus complexe et coûteuse en ressources. La recherche a pour but d'évaluer si une transcription détaillée, intégrant des éléments visuels, narratifs et sonores sous forme textuelle structurée, permettrait à un LLM de détecter efficacement des événements saillants dans une séquence vidéo. 

La contribution scientifique majeure de cette étude serait de démontrer que des systèmes purement textuels peuvent effectivement remplacer, dans certaines situations précises, des méthodes multimodales gourmandes en puissance de calcul, tout en conservant une capacité d'analyse satisfaisante. 

Sur le plan pratique, les retombées potentielles incluent la réduction significative des ressources nécessaires pour réaliser l'analyse automatisée de vidéos. Des mesures quantitatives seront présentées comme résultat, comparant les ressources utilisées par une intelligence artificielle multimodale à celles utilisées par notre modèle pour un résultat similaire de détection de moments forts. De telles avancées ouvriraient la voie à des robots autonomes capables de surveillance intelligente, d'inspection automatisée ou de reconnaissance rapide d'événements critiques dans des environnements variés (sécurité, exploration, assistance aux humains), tout en minimisant les contraintes matérielles et énergétiques. Elles permettraient également aux robots disposant de peu de ressources computationnelles de réaliser l’analyse vidéo directement sur leur matériel embarqué, renforçant la sécurité et la confidentialité des données en évitant leur transmission via internet et l’interaction avec des \textit{API} infonuagiques, un enjeu important pour des robots comme T-Top \cite{maheux2022_ttop}.

À plus long terme, ce projet contribuerait à la communauté scientifique en clarifiant les avantages et les limites d'une représentation purement textuelle pour des contenus audiovisuels. Il permettrait également de mieux cerner les conditions sous lesquelles l'information visuelle peut être traduite efficacement en descripteurs textuels sans perte critique de données, ainsi que les risques potentiels liés à une simplification excessive.


